\section{Why individual electrons have no good quantum numbers}\label{sec:qn}

\begin{reviewclaim}
``\dots according to modern quantum mechanics, in the stationary states the quantum numbers of individual
electrons do not commute with the Hamiltonian, only the total angular momentum does. For this reason
\dots\ the Pauli exclusion principle has been reformulated to reflect the antisymmetry of the
wavefunction \dots'' (p.~448, ref.~77 Scerri 1991). ``\dots assigning individual electrons to individual
orbitals and individual sets of quantum numbers is erroneous'' (p.~441).
\end{reviewclaim}

\subsection{The commutators}
Take an atom (nucleus at the origin), so $\Hop=\sum_i\hat h(i)+\sum_{i<j}r_{ij}^{-1}$ with
$\hat h(i)=-\tfrac12\nabla_i^2-Z/r_i$. Orbital angular momentum of electron 1:
$\hat{\mathbf l}_1=-i\,\mathbf r_1\times\nabla_1$.
\begin{enumerate}
  \item $\hat{\mathbf l}_1$ commutes with $\hat h(1)$ (spherical symmetry) and with $\hat h(j)$, $j\ne1$ (different variables).
        The only possible obstruction is $1/r_{12}$.
  \item Apply $\hat l_{1z}=-i(x_1\partial_{y_1}-y_1\partial_{x_1})$ to $V=1/r_{12}$, with
        $\partial_{y_1}r_{12}^{-1}=-(y_1-y_2)/r_{12}^3$ and $\partial_{x_1}r_{12}^{-1}=-(x_1-x_2)/r_{12}^3$:
        \[
          \hat l_{1z}V=-i\,\frac{-x_1(y_1-y_2)+y_1(x_1-x_2)}{r_{12}^3}=-i\,\frac{x_1y_2-x_2y_1}{r_{12}^3}.
        \]
        Since $\hat l_{1z}$ is a first-order differential operator, $[\hat l_{1z},V]f=(\hat l_{1z}V)f$ for any $f$, so
        \begin{equation}
          [\hat l_{1z},\,r_{12}^{-1}]=-i\,\frac{x_1y_2-x_2y_1}{r_{12}^3}\neq0 .
        \end{equation}
  \item For electron 2, $\partial_{y_2}r_{12}^{-1}=+(y_1-y_2)/r_{12}^3$ etc., which gives
        $\hat l_{2z}V=-i\,(x_2y_1-x_1y_2)/r_{12}^3=-\hat l_{1z}V$.
  \item Hence $[\hat L_z,r_{12}^{-1}]=[\hat l_{1z}+\hat l_{2z},r_{12}^{-1}]=0$. The same holds for the $x$ and $y$
        components. Physically, $\hat{\mathbf L}$ generates a simultaneous rotation of both electrons, which
        leaves $r_{12}$ unchanged; $\hat{\mathbf l}_1$ rotates electron 1 alone, which changes $r_{12}$.
  \item Consequently $[\Hop,\hat{\mathbf L}]=0$ and $[\Hop,\hat L^2]=0$, but $[\Hop,\hat{\mathbf l}_1]\neq0$ and $[\Hop,\hat l_1^2]\neq0$.
        $\hat{\mathbf S}$ and $\hat S^2$ commute with $\Hop$ trivially, since $\Hop$ is spin-free.
\end{enumerate}
\begin{resultbox}
Exact atomic eigenstates can be labelled by $L,M_L,S,M_S$ (and parity), not by one-electron $(n,l,m_l,m_s)$.
Those labels belong to the independent-particle model, in which $1/r_{ij}$ is replaced by a spherically
averaged mean field. Pauli's 1925 one-electron formulation is therefore a statement about that model. The
exact, model-independent statement is the antisymmetry \cref{eq:antisym}.
\end{resultbox}
\begin{verifybox}
T13 (\texttt{t10\_t12\_t13\_symbolic.py}): $[\hat l_{1z},r_{12}^{-1}]f$ evaluated symbolically on an arbitrary
function $f$ reproduces the expression above. $[\hat L_k,r_{12}^{-1}]f=0$ for $k=x,y,z$, and $[\hat l_1^2,r_{12}^{-1}]\neq0$. \textbf{PASS}.
\end{verifybox}
\begin{sourcebox}
Scerri, \emph{Br.\ J.\ Philos.\ Sci.}\ \textbf{42}, 309 (1991); Pauli, \emph{Z.\ Phys.}\ \textbf{31}, 765 (1925): not
accessed. The commutator algebra is standard and was verified here.
\end{sourcebox}
